\section*{Chapter \ref{ch:g11:oligopoly-and-games} --- Oligopoly and a First Look at Games}
% figdata: figdata/grade-11/05-oligopoly-and-games/out/ (opec, overcharge); ledger g11:opec-1980s, g11:lysine, g11:connor2005-*; Kestria numbers invented.

\begin{solution}{exo:g11:oligopoly-and-games:1}
Top is \omterm{def:g11:oligopoly-and-games:dominant-strategy}{dominant} for the row player ($4 > 3$ and $2 > 1$); Right is \omterm{def:g11:oligopoly-and-games:dominant-strategy}{dominant} for the column player ($3 > 1$ and $2 > 0$). The outcome is (Top, Right), with \omterm{def:g11:oligopoly-and-games:game}{payoffs} (2, 3).
\end{solution}

\begin{solution}{exo:g11:oligopoly-and-games:2}
Low is \omterm{def:g11:oligopoly-and-games:dominant-strategy}{dominant} for each café ($12 > 10$ and $6 > 4$). Both charge low and earn 6 each, although both high would give 10 each.
\end{solution}

\begin{solution}{exo:g11:oligopoly-and-games:3}
Top is \omterm{def:g11:oligopoly-and-games:dominant-strategy}{dominant} for the row player ($5 > 4$, $3 > 2$). The column player has none: Left is best against Top, Right against Bottom. Since the row player plays Top, the column
player plays Left: the equilibrium is (Top, Left), \omterm{def:g11:oligopoly-and-games:game}{payoffs} (5, 2).
\end{solution}

\begin{solution}{exo:g11:oligopoly-and-games:4}
Both A and both B: in each, a firm that switched alone would get 0. Both firms prefer the first.
\end{solution}

\begin{solution}{exo:g11:oligopoly-and-games:5}
Saudi output fell by $9.81 - 3.39 = 6.42$ million barrels a day (about 65\,\%); the others' rose from 11.42 to 11.92. In 1986 Saudi output rose to 4.87 and the others' to 12.60,
and the real price of petrol fell from 3.50 to 2.65 dollars a gallon.
\end{solution}

\begin{solution}{exo:g11:oligopoly-and-games:6}
Both 22.5: $p = 55$, each $45 \times 22.5 = 1012.5$. One at 30, the other at 22.5: $p = 47.5$, profits $37.5 \times 30 = 1125$ and $37.5 \times 22.5 = 843.75$. Both 30: $p = 40$,
each 900. Producing 30 is \omterm{def:g11:oligopoly-and-games:dominant-strategy}{dominant} ($1125 > 1012.5$, $900 > 843.75$); the equilibrium is (30, 30), a \omterm{def:g11:oligopoly-and-games:prisoners-dilemma}{prisoner's dilemma}.
\end{solution}

\begin{solution}{exo:g11:oligopoly-and-games:7}
Both 30: $p = 70$, each $60 \times 30 = 1800$. 40 against 30: $p = 60$, profits $50 \times 40 = 2000$ and $50 \times 30 = 1500$. Both 40: $p = 50$, each $40 \times 40 = 1600$. 40 is
\omterm{def:g11:oligopoly-and-games:dominant-strategy}{dominant}; the equilibrium is (40, 40), with 1600 each against 1800.
\end{solution}

\begin{solution}{exo:g11:oligopoly-and-games:8}
$\pi(q) = (130 - q - 30 - 10)q = (90 - q)q$, so $\pi'(q) = 90 - 2q = 0$ at $q = 45$; $p = 55$ and $\pi = 45 \times 45 = 2025$, a gain of $2025 - 1800 = 225$.
\end{solution}

\begin{solution}{exo:g11:oligopoly-and-games:9}
Against silence, silence is best ($60 > 0$); against reporting, reporting is best ($-100 > -200$). Both (silent, silent) and (report, report) are equilibria. A member that
thinks the other may report does better to report first: if the chance that the other reports is $r$, staying silent gives $60(1 - r) - 200r$ and reporting $-100r$, so
reporting is better as soon as $r > 0.375$. Mutual suspicion can push the \omterm{def:g11:oligopoly-and-games:cartel}{cartel} into the reporting equilibrium.
\end{solution}

\begin{solution}{exo:g11:oligopoly-and-games:10}
The median international overcharge is 32\,\%: without the \omterm{def:g11:oligopoly-and-games:cartel}{cartel} the goods would have cost $660/1.32 = 500$ million euros, so buyers paid 160 million euros more.
\end{solution}

\begin{solution}{exo:g11:oligopoly-and-games:11}
(Top, Left) and (Bottom, Right). The row player prefers the first (3 against 2), the column player the second.
\end{solution}

\begin{solution}{exo:g11:oligopoly-and-games:12}
The lysine \omterm{def:g11:oligopoly-and-games:cartel}{cartel} lasted for years and raised prices by 70\,\% within nine months: its members met again and again and watched each other, so the assumption that the \omterm{def:g11:oligopoly-and-games:game}{game}
is played once fails. In laboratory one-shot dilemmas many players cooperate: the assumption that each player cares only about her own money \omterm{def:g11:oligopoly-and-games:game}{payoff} fails for them.
\end{solution}

\begin{solution}{exo:g11:oligopoly-and-games:13}
Keeping to the quota gives $\pi_c + \pi_c\,\delta/(1 - \delta)$; cheating gives $\pi_d + \pi_n\,\delta/(1 - \delta)$. The first is at least the second when $\pi_d - \pi_c \leq (\pi_c -
\pi_n)\,\delta/(1 - \delta)$. Quarries: $144 \leq 128\,\delta/(1 - \delta)$ gives $144 \leq 272\,\delta$, $\delta \geq 144/272 \approx 0.53$.
\end{solution}

\begin{solution}{exo:g11:oligopoly-and-games:14}
New sellers enter while there are profits, taking buyers from each seller, until each seller's demand curve just touches its average cost curve: at that point the price
equals average cost, so the profit is zero. The demand curve slopes down, so it touches the average cost curve where that curve slopes down too, to the left of its lowest
point; and since \omterm{def:g11:marginal-reasoning:marginal-benefit}{marginal revenue} is below the price, so is \omterm{def:g11:marginal-reasoning:marginal-benefit}{marginal cost}.
\end{solution}

\begin{solution}{exo:g11:oligopoly-and-games:15}
Both left (1, 1); both right (1, 1); different sides ($-10$, $-10$). Both matching combinations are equilibria. Once a law names one side, each driver's best reply to the
others obeying it is to obey it too: the law only picks the equilibrium. In a \omterm{def:g11:oligopoly-and-games:prisoners-dilemma}{prisoner's dilemma} each player's best reply is to break an agreement, so a rule works only if a
penalty changes the \omterm{def:g11:oligopoly-and-games:game}{payoffs}.
\end{solution}

\begin{solution}{pb:g11:oligopoly-and-games:1}
\textbf{1.} $120 - 2Q = 24$ gives $Q = 48$; $p = 72$.

\textbf{2.} Quota 24 each; profit $(72 - 24) \times 24 = 1152$.

\textbf{3.} At $p = 24$, $Q = 96$; \omterm{def:g10:what-prices-do:dwl}{deadweight loss} $\tfrac12 \times (96 - 48) \times (72 - 24) = 1152$.

\textbf{4.} $\pi_A = (120 - q_A - 24 - 24)\,q_A = (72 - q_A)\,q_A$.

\textbf{5.} $72 - 2q_A = 0$ gives $q_A = 36$; $p = 120 - 60 = 60$; $\pi_A = 36 \times 36 = 1296$.

\textbf{6.} $\pi_A = (96 - q_A - q_B)\,q_A$ is largest where $96 - 2q_A - q_B = 0$: $q_A = (96 - q_B)/2$.

\textbf{7.} $q = (96 - q)/2$ gives $q = 32$ for each.

\textbf{8.} $p = 120 - 64 = 56$; each earns $32 \times 32 = 1024$.

\textbf{9.} Both 24: (1152, 1152). A 24, B 32: price 64, (960, 1280). A 32, B 24: (1280, 960). Both 32: (1024, 1024).

\textbf{10.} 32 is \omterm{def:g11:oligopoly-and-games:dominant-strategy}{dominant} for each ($1280 > 1152$, $1024 > 960$), and the equilibrium gives each 1024, less than the 1152 of keeping to the quota.

\textbf{11.} $1296 - 1152 = 144$.

\textbf{12.} At a price of 60, B earns $(60 - 24) \times 24 = 864$.

\textbf{13.} $1152 - 1024 = 128$ a year.

\textbf{14.} One year (128) is less than 144; two years (256) are more. Two years of punishment remove the gain.

\textbf{15.} $128/1.1 + 128/1.21 \approx 116.4 + 105.8 = 222.2 > 144$: yes. One year alone, 116.4, is not enough.

\textbf{16.} No: nothing can punish cheating in the last year, so each quarry plays its \omterm{def:g11:oligopoly-and-games:dominant-strategy}{dominant strategy} and produces 32.

\textbf{17.} In the ninth year each knows the tenth will be played at 32 whatever happens, so the threat of punishment is empty and each produces 32 in the ninth year too;
the same reasoning runs back to the first year, and the \omterm{def:g11:oligopoly-and-games:cartel}{cartel} never works.

\textbf{18.} If there is always a chance that the \omterm{def:g11:oligopoly-and-games:game}{game} goes on, there is always a future that a cheater can lose, and the trigger \omterm{def:g11:oligopoly-and-games:game}{strategy} keeps its force.

\textbf{19.} A member can now cheat and report, escaping the fine, while the other is fined: the punishment that holds the \omterm{def:g11:oligopoly-and-games:cartel}{cartel} together weighs less against the gains of
leaving it, and each member must also fear being reported.

\textbf{20.} Cheating on the \omterm{def:g11:oligopoly-and-games:cartel}{cartel}'s quota brings a quarry 144 thousand kests in the year it cheats; a punishment in which both produce 32 costs it 128 thousand kests a year,
so two years of punishment remove the gain (one year is not enough).
\end{solution}
